%%%%%
%
%% Dans ce fichier, les �ventuels accents sont cod�s
%% en ISO-8859-15 (aka latin9)
%% La ligne suivante donne trois exemples 
%% e accent aigu : �, e accent grave : �, e-dans-l'o : �
\section{Indice de stationnarit� de $U_*$}

\itemstyle{Op\'erateur~:} \operfmt{Stationnarite\_Ustar}
\index{Stationnarite\_Ustar}

%%%%%OBJET%%%%
\itemstyle{Objet~:}
Calcul d'un indice de stationnarit�
de la vitesse de frottement $U_*$ en comparant la moyenne de $U_*$ calcul�e sur des sous-p�riodes, et $U_*$ calcul�e sur la totalit� de la p�riode.

%%%%%ARGUMENTS%%%%
\bigskip

\begin{tabarg}
\multicolumn{3}{l}{\titretab {Entr\'ees}}\\
\hline
$U$ & Vecteur, \texttt{float} & 1�re composante du vent horizontal (m/s) \\
$V$ & Vecteur, \texttt{float} & 2�me composante du vent horizontal (m/s) \\
$W$ & Vecteur, \texttt{float} & Composante du vent vertical (m/s) \\
$f_e$ & Coeff., \texttt{float} & Fr�quence d'�chantillonnage X et Y (Hz)\\
$T_1$ & Coeff., \texttt{int} & Dur�e de la p�riode d'int�gration partielle  (s)\\
$T_2$ & Coeff., \texttt{int} & Dur�e de la p�riode d'int�gration totale  (s)\\
$R_{manq}$ & Coeff., \texttt{int} & Fraction minimale de valeurs non manquantes tol�r�es (compris entre 0 et 1) \\
\end{tabarg}

\medskip
%%%%%SORTIES%%%%
\begin{tabarg}
\multicolumn{3}{l}{\titretab {Sorties}}\\
\hline
$I_{stat}$ & Vecteur, \texttt{float} & Indice de stationnarit� (sans unit�)\\
\end{tabarg}


%%%%%FORMULE%%%%
\itemstyle{Formule}

%% Formule de l'operateur : Stationnarite_Ustar
%% Dans ce fichier, les �ventuels accents sont cod�s
%% en ISO-8859-15 (aka latin9)
%% La ligne suivante donne trois exemples 
%% e accent aigu : �, e accent grave : �, e-dans-l'o : �

$M$ est le nombre de sous-p�riodes sur lesquelles on calcule $U_*$. 
Donc $M$ vaut $T_2 /T_1$ ( $T_2$ > $T_1$ et $T_2$ multiple de $T_1$ )\\

$U_*$ sur la $i$�me p�riode $T_1$ est calcul�e par :
$$
(U_*)_i = \left( \left ( \overline{U'W'} \right)^2_i + 
                 \left ( \overline{V'W'} \right)^2_i \right)^{1/4} 
$$

Puis $U_*$ moyen vaut
$$
U_{*moy} = \frac{1}{M} \sum_i (U_*)_i
$$
L'indice de stationnarit� vaut :
$$
Istat = \left\vert  \frac{ U_*  - U_{*moy} }
                   { U_* } \right\vert
$$
avec $U_*$ : vitesse de frottement calcul�e sur toute la p�riode $T_2$

%%%%%AUTEUR%%%%
\itemstyle{Auteur~:}
CNRM/GMEI/TRAMM

%%%%%REFERENCE%%%%
\itemstyle{R\'ef\'erences~:}
Micro-meteorology, T. Foken, 2008 \cite{MICROMETEO-FOKEN}
